\documentclass[10pt,a4paper]{article}

% ---------- lingua, codifica, font (testo + matematica) ----------
\usepackage[T1]{fontenc}
\usepackage[utf8]{inputenc}
\usepackage[italian]{babel}
\usepackage{mathpazo}                 % Palatino per testo e matematica
\usepackage{microtype}

% ---------- matematica ----------
\usepackage{amsmath,amssymb,amsfonts,mathtools,bm,cancel,mathrsfs,nicefrac}
\DeclareMathOperator{\sen}{sen}       % seno all'italiana
\DeclareMathOperator{\tg}{tg}
\DeclareMathOperator{\ctg}{cotg}
\DeclareMathOperator{\arcsen}{arcsen}
\DeclareMathOperator{\arctg}{arctg}
\newcommand{\dd}{\mathrm{d}}                       % differenziale
\newcommand{\vv}[1]{\vec{#1}}                      % vettore con freccia
\newcommand{\norma}[1]{\left\lvert#1\right\rvert} % modulo di vettore
\newcommand{\pedice}[1]{_{\text{#1}}}              % pedice testuale: E\pedice{in}

% ---------- unità di misura e numeri (virgola decimale) ----------
\usepackage{siunitx}
\sisetup{output-decimal-marker={,},per-mode=symbol,
         inter-unit-product=\ensuremath{{}\cdot{}},group-digits=false}

% ---------- grafici e figure ----------
\usepackage{tikz}
\usetikzlibrary{arrows.meta,calc,positioning,patterns,decorations.pathmorphing}
\usepackage{pgfplots}
\pgfplotsset{compat=1.16}
\usepackage{graphicx}
% Grafici collocati nel punto pertinente dello svolgimento.
\newenvironment{grafico}{\begin{figure}[h!]\centering}{\end{figure}}

% ---------- impaginazione ----------
\usepackage[a4paper,left=1.8cm,right=1.8cm,top=2.3cm,bottom=2.2cm]{geometry}
\usepackage{fancyhdr,lastpage,needspace,enumitem,booktabs,array,colortbl,tabularx}
\setlength{\parindent}{0pt}
\setlength{\parskip}{3pt}
\setlength{\headheight}{14pt}
\setlist[enumerate]{leftmargin=1.7em,itemsep=2pt,topsep=2pt,parsep=0pt}
\setlist[itemize]{leftmargin=1.3em,itemsep=2pt,topsep=2pt,parsep=0pt}

% ---------- colori degli indicatori ----------
\usepackage[table,dvipsnames]{xcolor}
\definecolor{anac}{HTML}{C62828}\definecolor{anabg}{HTML}{FDECEA} % rosso
\definecolor{svic}{HTML}{1A56DB}\definecolor{svibg}{HTML}{E8F0FE} % blu
\definecolor{argc}{HTML}{A67C00}\definecolor{argbg}{HTML}{FFF6CC} % giallo
\definecolor{comc}{HTML}{2E7D32}\definecolor{combg}{HTML}{E6F4EA} % verde
% Solo griglie d'esame (su richiesta): IND arancio, DAT teal
\definecolor{indc}{HTML}{E65100}\definecolor{indbg}{HTML}{FFF0E0}
\definecolor{datc}{HTML}{00796B}\definecolor{datbg}{HTML}{E0F2F1}

% Riferimenti agli indicatori: usare queste macro in tutto il documento.
\newcommand{\coloreANA}{anac}\newcommand{\nomeANA}{Analizzare}
\newcommand{\coloreSVI}{svic}\newcommand{\nomeSVI}{Sviluppare e calcolare}
\newcommand{\coloreARG}{argc}\newcommand{\nomeARG}{Argomentare}
\newcommand{\coloreCOM}{comc}\newcommand{\nomeCOM}{Comunicare}
\newcommand{\coloreIND}{indc}\newcommand{\nomeIND}{Individuare}
\newcommand{\coloreDAT}{datc}\newcommand{\nomeDAT}{Interpretare i dati}
% Griglie d'esame: matematica \renewcommand{\nomeANA}{Comprendere}, \nomeSVI{Sviluppare il processo risolutivo};
% fisica/mista: \nomeSVI{Sviluppare il processo risolutivo}, \nomeDAT{Interpretare criticamente i dati} (mista: ...rappresentare, elaborare).
\newcommand{\sigla}[1]{\textcolor{\csname colore#1\endcsname}{\textbf{#1}}}
\newcommand{\indicatore}[1]{\textcolor{\csname colore#1\endcsname}{\textbf{#1: \csname nome#1\endcsname}}}

% ---------- box colorati e chip ----------
\usepackage{tcolorbox}
\tcbuselibrary{skins,breakable}

% Box passaggio chiave: \begin{boxana}[Titolo] ... \end{boxana}  (ana|svi|arg|com)
\newcommand{\defbox}[4]{%
  \newtcolorbox{box#1}[1][Passaggio chiave]{enhanced,breakable,frame hidden,
    borderline west={2.5pt}{0pt}{#2},colback=#3,boxrule=0pt,arc=1.5pt,
    left=7pt,right=5pt,top=2pt,bottom=2pt,before skip=4pt,after skip=4pt,
    fontupper=\small,before upper={\textbf{\idevidenza}\enspace\sigla{#4}\enspace\textbf{##1}\par}}}
\defbox{ana}{anac}{anabg}{ANA}
\defbox{svi}{svic}{svibg}{SVI}
\defbox{arg}{argc}{argbg}{ARG}
\defbox{com}{comc}{combg}{COM}
\defbox{ind}{indc}{indbg}{IND}
\defbox{dat}{datc}{datbg}{DAT}

% Box enunciato
\newtcolorbox{enun}{colback=black!3,colframe=black!25,boxrule=0.4pt,arc=2pt,
  left=6pt,right=6pt,top=3pt,bottom=3pt,fontupper=\small,
  before skip=5pt,after skip=5pt}

% Chip identificativo e massimo in punti; i dati sono forniti dal registro approvato.
\newcommand{\chip}[4]{\tcbox[on line,boxsep=0pt,left=3pt,right=3pt,top=1pt,bottom=1pt,
  arc=4pt,boxrule=0.4pt,colframe=#1,colback=#2,fontupper=\scriptsize\bfseries\color{#1}]{#3\ #4}}
\newcommand{\ANA}[2]{\chip{anac}{anabg}{\sigla{ANA}}{#1\, / massimo #2\,p}}
\newcommand{\SVI}[2]{\chip{svic}{svibg}{\sigla{SVI}}{#1\, / massimo #2\,p}}
\newcommand{\ARG}[2]{\chip{argc}{argbg}{\sigla{ARG}}{#1\, / massimo #2\,p}}
\newcommand{\COM}[2]{\chip{comc}{combg}{\sigla{COM}}{#1\, / massimo #2\,p}}
\newcommand{\IND}[2]{\chip{indc}{indbg}{\sigla{IND}}{#1\, / massimo #2\,p}}
\newcommand{\DAT}[2]{\chip{datc}{datbg}{\sigla{DAT}}{#1\, / massimo #2\,p}}
\newcommand{\ess}{\ensuremath{\bigstar}}
\newcommand{\legenda}{\indicatore{ANA}\quad\indicatore{SVI}\par
  \indicatore{ARG}\quad\indicatore{COM}}
% Griglie d'esame: \legendaesame{ANA}{IND}{SVI}{ARG} (matematica) o \legendaesame{ANA}{SVI}{DAT}{ARG} (fisica, mista)
\newcommand{\legendaesame}[4]{\indicatore{#1}\quad\indicatore{#2}\par\indicatore{#3}\quad\indicatore{#4}}
% Limite categorico sull'intera verifica; il registro non incrementa il conteggio.
\newcounter{numeroevidenze}
\newcommand{\controllaevidenze}{\stepcounter{numeroevidenze}%
  \ifnum\value{numeroevidenze}>40\relax
    \PackageError{correttore}{Superato il limite categorico di 40 evidenze}{Accorpare le evidenze prima di generare il correttore.}%
  \fi}
% Metadati mantenuti nel sorgente; nel corpo il titolo viene emesso solo dal box.
\newcommand{\idevidenza}{}
\newcommand{\evidenza}[5]{\controllaevidenze\def\idevidenza{#1}\needspace{8\baselineskip}}
% Registro raggruppato: esercizio, argomento comune, righe ID/nome/punto/indicatore/massimo.
\newcommand{\registrogruppo}[3]{\needspace{6\baselineskip}\par\smallskip{\small
  \textbf{Esercizio #1}\enspace #2\par
  \begin{tabularx}{\linewidth}{@{}lXllr@{}}\toprule
  ID & Evidenza & Punto & Ind. & Max\\\midrule
  #3\bottomrule\end{tabularx}\par}}
\newcommand{\inizioesercizi}{\clearpage}
\newcommand{\riepiloghi}[2]{\sez{Riepiloghi}
  \noindent\begin{minipage}[t]{0.58\linewidth}#1\end{minipage}\hfill
  \begin{minipage}[t]{0.39\linewidth}#2\end{minipage}\par}

\newcommand{\livelli}[5]{\begin{description}[leftmargin=2.7cm,style=nextline,font=\small\bfseries]
\item[Non raggiunto] #1
\if\relax\detokenize{#2}\relax\else\item[Iniziale] #2\fi
\item[Base] #3
\if\relax\detokenize{#4}\relax\else\item[Intermedio] #4\fi
\item[Avanzato] #5
\end{description}}

% Intestazioni di colonna colorate
\newcommand{\hANA}{\cellcolor{anabg}\sigla{ANA}}
\newcommand{\hSVI}{\cellcolor{svibg}\sigla{SVI}}
\newcommand{\hARG}{\cellcolor{argbg}\sigla{ARG}}
\newcommand{\hCOM}{\cellcolor{combg}\sigla{COM}}
\newcommand{\hIND}{\cellcolor{indbg}\sigla{IND}}
\newcommand{\hDAT}{\cellcolor{datbg}\sigla{DAT}}

% ---------- macro di struttura ----------
% Sezione semplice
\newcommand{\sez}[1]{\needspace{4\baselineskip}\par\vspace{1em}%
  {\large\bfseries #1}\par\vspace{-0.4em}\rule{\linewidth}{0.6pt}\par\vspace{0.2em}}
% Esercizio: \esercizio{Titolo}{ANA}{SVI}{ARG}{COM}{TOT}
\newcommand{\esercizio}[6]{\needspace{8\baselineskip}\par\vspace{1em}%
  {\large\bfseries #1}\par\vspace{0.25em}}
% Punto: \punto{a) titolo}
\newcommand{\punto}[1]{\needspace{6\baselineskip}\par\vspace{0.7em}{\bfseries #1}\par\vspace{0.2em}}
% Tabella punteggi del punto: \tabpunti{ANA}{SVI}{ARG}{COM}{TOT}  (usare 0 se assente)
\newcommand{\tabpunti}[5]{\par\noindent\renewcommand{\arraystretch}{1.15}\small
  \begin{tabular}{*{5}{|>{\centering\arraybackslash}p{1.4cm}}|}\hline
  \hANA & \hSVI & \hARG & \hCOM & \textbf{Tot}\\\hline
  #1 & #2 & #3 & #4 & \textbf{#5}\\\hline\end{tabular}\par\smallskip}
% Griglie d'esame: \tabpuntiesame{ANA}{IND}{SVI}{ARG}{v1}{v2}{v3}{v4}{TOT} (sigle nell'ordine dei valori)
\newcommand{\tabpuntiesame}[9]{\par\noindent\renewcommand{\arraystretch}{1.15}\small
  \begin{tabular}{*{5}{|>{\centering\arraybackslash}p{1.4cm}}|}\hline
  \csname h#1\endcsname & \csname h#2\endcsname & \csname h#3\endcsname & \csname h#4\endcsname & \textbf{Tot}\\\hline
  #5 & #6 & #7 & #8 & \textbf{#9}\\\hline\end{tabular}\par\smallskip}
% Nota (alternative, errori trascinati, criteri)
\newcommand{\nota}[1]{\begin{itemize}[font=\footnotesize]\item {\footnotesize #1}\end{itemize}}
\newcommand{\criteri}[1]{\par\nopagebreak[3]\smallskip\needspace{3\baselineskip}%
  {\footnotesize\textbf{Alternative, errori e criteri}\par
  \begin{itemize}[itemsep=1pt,topsep=1pt]#1\end{itemize}}}

% ---------- intestazione e piè di pagina ----------
\newcommand{\materia}{}
\newcommand{\classe}{}
\newcommand{\argomenti}{}
\pagestyle{fancy}\fancyhf{}
\renewcommand{\headrulewidth}{0.3pt}
\lhead{\small\textsc{Correttore\ \materia}}
\rhead{}
\cfoot{\small Pag.\ \thepage\ di \pageref{LastPage}}

\renewcommand{\materia}{Matematica}
\renewcommand{\classe}{Classe terza, liceo scientifico \textperiodcentered\ Rette e fasci \textperiodcentered\ Versione A}
\begin{document}
\begin{center}
{\LARGE\bfseries Correttore\ \materia}\\[3pt]
\classe
\end{center}

\legenda
\sez{Registro delle evidenze}

\registrogruppo{1}{RET, DIS, PAR (associazione indicata per ogni evidenza)}{
E01 & Imposta la retta per due punti, RET & a & \sigla{ANA} & 40\\
E02 & Ricava l'equazione della retta $AB$, RET & a & \sigla{SVI} & 60\\
E03 & Calcola $|AB|$, DIS & a & \sigla{SVI} & 40\\
E04 & Riconosce l'altezza come distanza punto-retta, DIS & b & \sigla{ANA} & 50\\
E05 & Calcola la distanza di $C$ da $AB$, DIS & b & \sigla{SVI} & 70\\
E06 & Calcola l'area, DIS & b & \sigla{SVI} & 40\\
E07 & Conferma l'area con un secondo metodo, DIS & b & \sigla{ARG} & 30\\
E08 & Imposta la parallela per $C$, PAR & c & \sigla{ANA} & 40\\
E09 & Scrive la parallela e la interseca con l'asse $x$, PAR & c & \sigla{SVI} & 60\\}
\registrogruppo{2}{PAR}{
E10 & Traduce la perpendicolarità in una condizione sui coefficienti angolari & intero & \sigla{ANA} & 50\\
E11 & Risolve per $k$ ed esamina la retta verticale & intero & \sigla{SVI} & 70\\}
\registrogruppo{3}{PAR}{
E12 & Individua punto medio e direzione dell'asse & intero & \sigla{ANA} & 50\\
E13 & Scrive l'equazione dell'asse & intero & \sigla{SVI} & 100\\
E14 & Verifica l'equidistanza di un punto dell'asse & intero & \sigla{ARG} & 50\\}
\registrogruppo{4}{FAS}{
E15 & Riconosce generatrici e centro del fascio & intero & \sigla{ANA} & 60\\
E16 & Determina centro, retta per il punto e parallela all'asse $x$ & intero & \sigla{SVI} & 80\\
E17 & Spiega perché la prima generatrice non si ottiene per alcun $k$ & intero & \sigla{ARG} & 50\\}
\registrogruppo{Intero elaborato}{TRASVERSALE, Competenze trasversali}{
E18 & Usa notazione e passaggi ordinati e leggibili & Intero elaborato & \sigla{COM} & 60\\}

\riepiloghi{%
\begin{tabular}{@{}lrrrrr@{}}\toprule
 & \hANA & \hSVI & \hARG & \hCOM & Tot\\\midrule
Esercizio 1 & 130 & 210 & 70 & 0 & 430\\
Esercizio 2 & 50 & 70 & 0 & 0 & 120\\
Esercizio 3 & 50 & 100 & 50 & 0 & 200\\
Esercizio 4 & 60 & 80 & 50 & 0 & 190\\
Intero elaborato & 0 & 0 & 0 & 60 & 60\\\midrule
Totale & 290 & 520 & 130 & 60 & 1000\\\bottomrule
\end{tabular}}{%
\begin{tabular}{@{}lr@{}}\toprule
Argomento & Max\\\midrule
RET, Equazione della retta & 100\\
DIS, Distanze & 230\\
PAR, Parallelismo e luoghi & 420\\
FAS, Fasci di rette & 190\\
Trasversale & 60\\\midrule
Totale & 1000\\\bottomrule
\end{tabular}}

\inizioesercizi
\esercizio{Esercizio 1 -- Triangolo $ABC$}{}{}{}{}{}
\begin{enun}
Dati $A(-2,1)$, $B(4,3)$, $C(2,-3)$: (a) retta $AB$ e lunghezza di $AB$; (b) distanza di $C$ da $AB$, area di $ABC$ e verifica con un secondo metodo; (c) parallela ad $AB$ per $C$ e sua intersezione con l'asse $x$.
\end{enun}

\punto{a) Retta $AB$ e lunghezza}
\evidenza{E01}{1/a}{ANA}{RET}{40}
\begin{boxana}[Imposta la retta per due punti]
\begin{enumerate}
\item Poiché $x_A\neq x_B$ la retta non è verticale e si scrive $y-y_A=m(x-x_A)$.
\item Il coefficiente angolare si ricava dai dati: $m=\dfrac{y_B-y_A}{x_B-x_A}$.
\end{enumerate}
\end{boxana}
\evidenza{E02}{1/a}{SVI}{RET}{60}
\begin{boxsvi}[Ricava l'equazione della retta $AB$]
\begin{enumerate}
\item Sostituendo i dati: $m=\dfrac{3-1}{4-(-2)}=\dfrac13$.
\item $y-1=\dfrac13(x+2)$.
\item In forma implicita: $x-3y+5=0$.
\end{enumerate}
\criteri{\item Ammesse la retta per due punti in forma di determinante o il sistema con $y=mx+q$.
\item Errore su $m$ penalizzato qui; le richieste successive che lo usano coerentemente conservano il punteggio.}
\end{boxsvi}
\evidenza{E03}{1/a}{SVI}{DIS}{40}
\begin{boxsvi}[Calcola $|AB|$]
\begin{enumerate}
\item $|AB|=\sqrt{(x_B-x_A)^2+(y_B-y_A)^2}$.
\item $|AB|=\sqrt{6^2+2^2}=\sqrt{40}=2\sqrt{10}$.
\end{enumerate}
\end{boxsvi}

\punto{b) Distanza, area e verifica}
\evidenza{E04}{1/b}{ANA}{DIS}{50}
\begin{boxana}[Riconosce l'altezza come distanza punto-retta]
\begin{enumerate}
\item L'altezza relativa al lato $AB$ è la distanza di $C$ dalla retta $AB$, misurata sulla perpendicolare.
\item Si usa $d=\dfrac{|ax_0+by_0+c|}{\sqrt{a^2+b^2}}$ con la retta in forma implicita.
\end{enumerate}
\end{boxana}
\evidenza{E05}{1/b}{SVI}{DIS}{70}
\begin{boxsvi}[Calcola la distanza di $C$ da $AB$]
\begin{enumerate}
\item Retta $AB$: $x-3y+5=0$, punto $C(2,-3)$.
\item $h=\dfrac{|2-3(-3)+5|}{\sqrt{1^2+(-3)^2}}=\dfrac{16}{\sqrt{10}}=\dfrac{8\sqrt{10}}{5}$.
\end{enumerate}
\criteri{\item Nota del docente: ammesso anche il risultato approssimato $5{,}06$, purché la forma esatta compaia in un passaggio.}
\end{boxsvi}
\evidenza{E06}{1/b}{SVI}{DIS}{40}
\begin{boxsvi}[Calcola l'area del triangolo]
\begin{enumerate}
\item $\mathcal{A}=\dfrac12\,|AB|\,h$.
\item $\mathcal{A}=\dfrac12\cdot2\sqrt{10}\cdot\dfrac{16}{\sqrt{10}}=16$.
\end{enumerate}
\end{boxsvi}
\evidenza{E07}{1/b}{ARG}{DIS}{30}
\begin{boxarg}[Conferma l'area con un secondo metodo]
\begin{enumerate}
\item Con i vertici: $\mathcal{A}=\dfrac12\,\bigl|x_A(y_B-y_C)+x_B(y_C-y_A)+x_C(y_A-y_B)\bigr|$.
\item $\mathcal{A}=\dfrac12\,|-2\cdot6+4\cdot(-4)+2\cdot(-2)|=\dfrac12\cdot32=16$.
\item I due metodi concordano: il risultato è coerente.
\end{enumerate}
\end{boxarg}

\punto{c) Parallela per $C$}
\evidenza{E08}{1/c}{ANA}{PAR}{40}
\begin{boxana}[Imposta la parallela ad $AB$ per $C$]
\begin{enumerate}
\item Rette parallele non verticali hanno lo stesso coefficiente angolare: $m=\dfrac13$.
\item La retta passa per $C(2,-3)$.
\end{enumerate}
\end{boxana}
\evidenza{E09}{1/c}{SVI}{PAR}{60}
\begin{boxsvi}[Scrive la parallela e la interseca con l'asse $x$]
\begin{enumerate}
\item $y+3=\dfrac13(x-2)$, cioè $x-3y-11=0$.
\item Con $y=0$ si ottiene $x=11$: il punto è $(11,0)$.
\end{enumerate}
\end{boxsvi}

\esercizio{Esercizio 2 -- Perpendicolarità con parametro}{}{}{}{}{}
\begin{enun}
Date $r:\ y=3x+2$ e $s:\ kx+(k+1)y-5=0$, determinare $k$ perché $s\perp r$ e discutere il caso senza coefficiente angolare.
\end{enun}
\evidenza{E10}{2}{ANA}{PAR}{50}
\begin{boxana}[Traduce la perpendicolarità sui coefficienti angolari]
\begin{enumerate}
\item $m_r=3$.
\item Per $s\perp r$ con $s$ non verticale: $m_s=-\dfrac1{m_r}=-\dfrac13$, con $m_s=-\dfrac{k}{k+1}$ e $k\neq-1$.
\end{enumerate}
\end{boxana}
\evidenza{E11}{2}{SVI}{PAR}{70}
\begin{boxsvi}[Risolve per $k$ ed esamina la retta verticale]
\begin{enumerate}
\item $-\dfrac{k}{k+1}=-\dfrac13\ \Rightarrow\ 3k=k+1\ \Rightarrow\ k=\dfrac12$.
\item Per $k=-1$ la retta è $-x-5=0$, cioè verticale $x=-5$: non è perpendicolare a $r$, che non è orizzontale.
\item L'unico valore è $k=\dfrac12$.
\end{enumerate}
\end{boxsvi}

\esercizio{Esercizio 3 -- Asse di un segmento}{}{}{}{}{}
\begin{enun}
Asse del segmento $PQ$ con $P(-1,4)$, $Q(5,0)$ e verifica dell'equidistanza di un suo punto.
\end{enun}
\evidenza{E12}{3}{ANA}{PAR}{50}
\begin{boxana}[Individua punto medio e direzione dell'asse]
\begin{enumerate}
\item L'asse è la perpendicolare a $PQ$ per il punto medio $M$ di $PQ$.
\item $M=\left(\dfrac{x_P+x_Q}{2},\dfrac{y_P+y_Q}{2}\right)$ e $m_{\text{asse}}=-\dfrac{1}{m_{PQ}}$.
\end{enumerate}
\end{boxana}
\evidenza{E13}{3}{SVI}{PAR}{100}
\begin{boxsvi}[Scrive l'equazione dell'asse]
\begin{enumerate}
\item $M=(2,2)$.
\item $m_{PQ}=\dfrac{0-4}{5+1}=-\dfrac23$, quindi $m_{\text{asse}}=\dfrac32$.
\item $y-2=\dfrac32(x-2)$, cioè $3x-2y-2=0$.
\end{enumerate}
\criteri{\item Ammesso il luogo dei punti equidistanti: $|XP|^2=|XQ|^2$ porta alla stessa equazione.}
\end{boxsvi}
\evidenza{E14}{3}{ARG}{PAR}{50}
\begin{boxarg}[Verifica l'equidistanza di un punto dell'asse]
\begin{enumerate}
\item Il punto $X(0,-1)$ soddisfa $3\cdot0-2\cdot(-1)-2=0$: sta sull'asse.
\item $|XP|^2=1+25=26$ e $|XQ|^2=25+1=26$: $X$ è equidistante da $P$ e da $Q$.
\end{enumerate}
\end{boxarg}

\esercizio{Esercizio 4 -- Fascio di rette}{}{}{}{}{}
\begin{enun}
Fascio $k(2x+y+1)+(-x+2y-3)=0$: centro, retta per $(3,0)$, parallela all'asse $x$ e motivazione sulla prima generatrice.
\end{enun}
\evidenza{E15}{4}{ANA}{FAS}{60}
\begin{boxana}[Riconosce generatrici e centro del fascio]
\begin{enumerate}
\item L'equazione è $k\,r_1+r_2=0$ con $r_1:\ 2x+y+1=0$ e $r_2:\ -x+2y-3=0$.
\item Il centro è il punto comune a $r_1$ e $r_2$: tutte le rette del fascio vi passano.
\end{enumerate}
\end{boxana}
\evidenza{E16}{4}{SVI}{FAS}{80}
\begin{boxsvi}[Determina centro, retta per il punto e parallela all'asse $x$]
\begin{enumerate}
\item Sistema $\begin{cases}2x+y+1=0\\-x+2y-3=0\end{cases}$: $y=-2x-1$ e $-5x-5=0$, quindi $C_0=(-1,1)$.
\item Per $(3,0)$: $7k-6=0$, $k=\dfrac67$, da cui $x+4y-3=0$.
\item Parallela all'asse $x$: il coefficiente di $x$, $2k-1$, è nullo, $k=\dfrac12$; si ottiene $y=1$.
\end{enumerate}
\end{boxsvi}
\evidenza{E17}{4}{ARG}{FAS}{50}
\begin{boxarg}[Spiega perché la prima generatrice non si ottiene]
\begin{enumerate}
\item $k\,r_1+r_2=0$ coincide con $r_1=0$ solo se $r_2$ è multipla di $r_1$.
\item $r_1$ e $r_2$ sono incidenti, quindi non proporzionali: nessun $k$ reale dà $r_1$ (si avrebbe solo per $k\to\infty$).
\end{enumerate}
\end{boxarg}

\esercizio{Intero elaborato}{}{}{}{}{}
\evidenza{E18}{Intero elaborato}{COM}{TRASVERSALE}{60}
\begin{boxcom}[Usa notazione e passaggi ordinati e leggibili]
\begin{enumerate}
\item Le equazioni sono scritte per intero con il segno $=$, gli estremi dei segmenti e i punti sono indicati.
\item I passaggi sono ordinati e ogni risultato è evidenziato.
\end{enumerate}
\criteri{\item Lo stesso errore di notazione ripetuto vale una volta sola. Non si valuta la grafia.}
\end{boxcom}
\end{document}
